\chapter{Modell}\label{sec:modell}

Tight-Binding Hamiltonian mit Hubbard-Anziehung gemäß gängiger Supraleitungstheorie

\begin{align}
    \mathcal{H} = \sum_{\vec{k}, \sigma} f(\vec{k}) A_{\vec{k}, \sigma}^{\dagger} B_{\vec{k}, \sigma} + h.c.  - U \sum_{\vec{R}} \sum_C C_{\vec{R}, \uparrow}^\dagger C_{\vec{R}, \uparrow} C_{\vec{R}, \downarrow}^\dagger C_{\vec{R}, \downarrow} \, ,
\end{align}

wobei 
\begin{equation}
    f(\vec{k}) = -t \sum_{\vec{\delta}} e^{i\vec{k}\cdot\vec{\delta}} = -t(e^{-ik_x a} + 2e^{ik_x \frac{a}{2}} \cos\left(\frac{k_y a \sqrt{3}}{2}\right))
\end{equation}
ist und $U>0$,$t>0$ die Hubbard-Anziehung und Hüpfparameter sind. Die Summe über die Operatoren $C$ soll hier bedeuten, dass über $A$ und $B$-Type Operatoren summiert wird.\\

Um die Tranlsationsinvarianz des Systems auszunutzen, wird der Hubbard-Term hier Fouriertransformiert.

\section{Fouriertrafo}

Fouriertrafo wird mit 
\begin{align}
    C_{\vec{R}} = \frac{1}{\sqrt{N}} \sum_{\vec{k}} C_{\vec{k}} e^{i \vec{k} \cdot \vec{R}}
\end{align}
angesetzt. \\

Damit gilt für den Wechselwirkungsteil des Hamiltonians:
\begin{align}
    \mathcal{H}_\text{WW} &= -U \sum_C \sum_{\vec{m}\vec{n}\vec{o}\vec{p}} C_{\vec{m}, \uparrow}^\dagger C_{\vec{n}, \uparrow} C_{\vec{o}, \downarrow}^\dagger C_{\vec{p}, \downarrow} \underbrace{\frac{1}{N^2} \sum_{\vec{R}} e^{i(-\vec{m}-\vec{o} + \vec{n} + \vec{p} )\cdot \vec{R}}}_{= \frac{1}{N} \delta_{\vec{m}+\vec{o}, \vec{n} + \vec{p}}} 
\end{align}
wobei durch $\delta_{\vec{m}+\vec{o}, \vec{n} + \vec{p}}$ die Impulserhaltung gewährleistet ist. Durch diese Bedingung kann die Anzahl der freien Parameter auf drei reduziert werden, sodass der Hamiltonian dann in die Form
\begin{align}
    \mathcal{H}_\text{WW} &= -U \sum_C \sum_{\vec{k}\, \vec{k'} \vec{q} } C_{\vec{k}+\vec{q}, \uparrow}^\dagger C_{\vec{k}, \uparrow} C_{\vec{k'} - \vec{q}, \downarrow}^\dagger C_{\vec{k'}, \downarrow}
\end{align}
gebracht werden kann. Um nun die Form eines typischen BCS-Hamiltonians zu erhlangen, können noch die Vertauschungsrelationen $\{C_{\vec{k}, \uparrow} ,C_{\vec{k'} - \vec{q}, \downarrow}^\dagger\} = 0$ und $\{C_{\vec{k}, \uparrow}, C_{\vec{k'}, \downarrow}\} = 0$ ausgenutzt werden. Damit ergibt sich
\begin{align}
    \mathcal{H}_\text{WW} &= -U \sum_C \sum_{\vec{k}\, \vec{k'} \vec{q} } C_{\vec{k}+\vec{q}, \uparrow}^\dagger C_{\vec{k'} - \vec{q}, \downarrow}^\dagger C_{\vec{k'}, \downarrow} C_{\vec{k}, \uparrow} \, . \label{eq:H_ww}
\end{align}

\section{Mean-Field-Näherung}

Um sich nicht mit quartischen Termen herumschlagen zu müssen, wird auf die Operatoren in Gleichung~\eqref{eq:H_ww} Wicks Theorem angewandt:
\begin{align*}
    C_{\vec{k}+\vec{q}, \uparrow}^\dagger C_{\vec{k'} - \vec{q}, \downarrow}^\dagger C_{\vec{k'}, \downarrow} C_{\vec{k}, \uparrow} &= \, :C_{\vec{k}+\vec{q}, \uparrow}^\dagger C_{\vec{k'} - \vec{q}, \downarrow}^\dagger C_{\vec{k'}, \downarrow} C_{\vec{k}, \uparrow}: \\
    &+ \, : C_{\vec{k}+\vec{q}, \uparrow}^\dagger C_{\vec{k'} - \vec{q}, \downarrow}^\dagger : \expval{C_{\vec{k'}, \downarrow} C_{\vec{k}, \uparrow}} + \expval{ C_{\vec{k}+\vec{q}, \uparrow}^\dagger  C_{\vec{k'} - \vec{q}, \downarrow}^\dagger } : C_{\vec{k'}, \downarrow} C_{\vec{k}, \uparrow} :\\
    &- : C_{\vec{k}+\vec{q}, \uparrow}^\dagger  C_{\vec{k'}, \downarrow}  :  \expval{ C_{\vec{k'} - \vec{q}, \downarrow}^\dagger  C_{\vec{k}, \uparrow} } -  \expval{ C_{\vec{k}+\vec{q}, \uparrow}^\dagger  C_{\vec{k'}, \downarrow} } :C_{\vec{k'} - \vec{q}, \downarrow}^\dagger  C_{\vec{k}, \uparrow}:  \\
    &+ \, : C_{\vec{k}+\vec{q}, \uparrow}^\dagger C_{\vec{k}, \uparrow} : \expval{ C_{\vec{k'} - \vec{q}, \downarrow}^\dagger C_{\vec{k'}, \downarrow} } + \expval{  C_{\vec{k}+\vec{q}, \uparrow}^\dagger C_{\vec{k}, \uparrow}  } : C_{\vec{k'} - \vec{q}, \downarrow}^\dagger C_{\vec{k'}, \downarrow}  : \\
    &+ \text{const.}
\end{align*}

Für die Mean-Field Näherung wird der quartische Term und auch die konstanten Terme vernachlässigt, da diese für die Dynamik des Systems keinen Unterschied machen. \\
Für die anderen drei Terme können die Erwartungswerte genauer untersucht werden.

Der erste Term enthält den anomalen Erwartungswert 
\begin{equation}
    \expval{C_{\vec{k'}, \downarrow} C_{\vec{k}, \uparrow}} = \delta_{\vec{k},\vec{-k'}}
\end{equation}
aufgrund der Impulserhaltung. Allgemein ist der Erwartungswert in der normalleitenden Phase null, in der supraleitenden Phase allerdings endlich. Er ist also später für den Ordnungsparameter der supraleitenden Phase relevant.

Für den negativen, zweiten Term gilt
\begin{equation}
    \expval{ C_{\vec{k'} - \vec{q}, \downarrow}^\dagger  C_{\vec{k}, \uparrow} } = 0
\end{equation}
aufgrund der Drehimpulserhaltung.

Der dritte Term hat Hartree-mäßige Erwartungswerte 
\begin{equation}
    \expval{ C_{\vec{k'} - \vec{q}, \downarrow}^\dagger C_{\vec{k'}, \downarrow} } = \delta_{\vec{q},0} \; .
\end{equation}

Damit kann der Wechselwirkungsterm also zwei verschiedenen Terme umgeschrieben werden. Der erste Term ergibt sich zu
\begin{align}
    \mathcal{H}_1 = -U \sum_C \sum_{\vec{k},\vec{q}} C_{\vec{k}+\vec{q}, \uparrow}^\dagger C_{-(\vec{k}+\vec{q}),\downarrow} \expval{ C_{-\vec{k}, \downarrow} C_{\vec{k}, \uparrow} } + \expval{ C_{\vec{k}+\vec{q}, \uparrow}^\dagger C_{-(\vec{k}+\vec{q}),\downarrow} }  C_{-\vec{k}, \downarrow} C_{\vec{k}, \uparrow} \; .
\end{align}

Dieser kann mit der Reparametrisierung $\vec{k'} = \vec{k} + \vec{q}$ weiter zu 
\begin{align}
    \mathcal{H}_1 = - \sum_C \Delta_C 
\end{align}

und 
\begin{align}
    
\end{align}
aufgeteit werden.